Hybrid vector search engines expose many physical plans---lexical retrieval, dense retrieval, learned-sparse retrieval, late interaction, fusion, metadata filtering, and reranking. Selecting a plan under a latency SLO is a query-optimization problem, yet static policies and per-query routers give little assurance that the chosen plan stays feasible once the next workload window shifts. We propose \CWC{}, a certified workload-window compiler for hybrid-search plan selection. Given a labeled historical workload, an unlabeled incoming workload window, and telemetry from candidate operators, \CWC{} partitions queries into workload cells, binds each cell to a physical plan, and abstains to a feasibility-best fallback when the selected plan is fragile. Its certificate is a finite-sample, distribution-free upper bound on the evaluation-window SLO-violation rate under a specified within-cell shift budget; in experiments, a conservative plug-in budget upper-bounds the realized violation in every run. On a $45$-block study over SciFact, FiQA, and NFCorpus with BM25, E5, Qwen3-Embedding, SPLADE, ColBERTv2, fusion, and reranking actions, \CWC{} ranks first on constrained utility, significantly outperforms four strong static operators, and statistically ties calibration-aware selectors while adding a deployment certificate they lack. At the binding budget it cuts realized SLO violations from about $11\%$ to about $0.1\%$ at roughly $4\%$ nDCG cost, and a scaling study shows the certified bound tightens with per-cell calibration size, as the theory predicts.
